\chapter{Intento Fallido: Recorte de Apollo}
\label{apendice:recorte_fallido}

Este apéndice documenta el primer enfoque que se intentó antes de decidir reimplementar desde cero: intentar recortar el módulo TLR del sistema Apollo completo.

\section{Estrategia Inicial}

La idea era aislar el módulo TLR del código completo de Apollo:

\begin{enumerate}
    \item Clonar el repositorio completo de Apollo
    \item Identificar el módulo de percepción de semáforos
    \item Eliminar progresivamente dependencias innecesarias
    \item Mantener solo el código esencial para TLR
\end{enumerate}

\section{Problemas Encontrados}

\subsection{Complejidad del Sistema}

\begin{itemize}
    \item \textbf{Tamaño del código:} Millones de líneas de código C++
    \item \textbf{Dependencias profundas:} El módulo TLR dependía de:
    \begin{itemize}
        \item CyberRT: Framework de mensajería propio de Apollo
        \item Librerías de utilidades internas
        \item Configuraciones específicas del ecosistema Apollo
    \end{itemize}
\end{itemize}

\subsection{Requisitos de Configuración}

Apollo requiere una configuración extensa:

\begin{table}[H]
\centering
\begin{tabular}{|l|l|}
\hline
\textbf{Componente} & \textbf{Requisito} \\
\hline
CPU & 8+ cores \\
RAM & 16 GB \\
GPU & NVIDIA RTX 20xx+ \\
OS & Ubuntu 18.04/20.04/22.04 \\
CUDA & 11.8 \\
Docker & $\geq$ 19.03 \\
Espacio disco & $\sim$50 GB \\
\hline
\end{tabular}
\caption{Requisitos de hardware/software de Apollo}
\end{table}

\subsection{Proceso de Compilación}

La compilación de Apollo involucra:
\begin{enumerate}
    \item Descargar imagen Docker ($\sim$15 GB)
    \item Configurar container con NVIDIA runtime
    \item Compilar todo el sistema (30+ minutos en hardware típico)
    \item Resolver dependencias de Bazel (sistema de build)
\end{enumerate}

\subsection{Errores Encontrados}

Durante los intentos de recorte, se encontraron múltiples errores:

\begin{itemize}
    \item Errores crípticos de Bazel al remover dependencias
    \item Conflictos de versiones entre librerías internas
    \item Problemas de linking al intentar compilar módulos aislados
    \item Errores de runtime por falta de configuraciones de CyberRT
\end{itemize}

\section{Tiempo Invertido}

Se invirtieron varias semanas intentando:
\begin{itemize}
    \item Configurar el ambiente de desarrollo
    \item Entender el sistema de build
    \item Identificar y remover dependencias
    \item Resolver errores de compilación
\end{itemize}

Sin lograr una versión ejecutable del módulo TLR aislado.

\section{Razones del Fracaso}

\begin{enumerate}
    \item \textbf{Acoplamiento estrecho:} El código de Apollo está diseñado para funcionar como un sistema integrado, no como módulos independientes

    \item \textbf{CyberRT como dependencia central:} El framework de mensajería está entrelazado con toda la lógica de negocio

    \item \textbf{Falta de documentación:} No hay guías para extraer módulos individuales

    \item \textbf{Complejidad del build:} Bazel hace difícil entender qué dependencias son realmente necesarias
\end{enumerate}

\section{Lección Aprendida}

Esta experiencia motivó el cambio de estrategia hacia la reimplementación. La lección principal fue:

\textbf{Para sistemas legacy complejos con alto acoplamiento, reimplementar un subsistema desde cero puede ser más viable que intentar extraerlo del sistema original.}

La reimplementación, aunque requiere más trabajo inicial de análisis, produce un artefacto mucho más manejable y comprensible.

\section{Valor del Intento Fallido}

Aunque el recorte no funcionó, el intento no fue tiempo perdido:

\begin{itemize}
    \item Reveló la complejidad real del sistema
    \item Motivó la búsqueda de alternativas
    \item Generó familiaridad inicial con el código de Apollo
    \item Informó las decisiones de diseño de la reimplementación
\end{itemize}